\shorthandoff{<>}% Babel español trata como abreviaturas los delimitadores de curvatura de Xy-pic.

\setcounter{subsection}{1}
\setcounter{subsubsection}{0}
\begin{sublemme}
\label{XIII.3.1.1}
Sean $S$ un esquema localmente noetheriano y
$$
\xymatrix@C=.5cm{ \Phi \ar[r]^p & \Phi_1 \ar@<0.6ex>[rr]^{p_1}
\ar@<-0.6ex>[rr]_{p_2} && \Phi_2 }
$$
un diagrama exacto de stacks sobre~$S$ \eqref{XIII.1.10.1}. Si $\Phi_1$
es constructible, lo mismo ocurre con~$\Phi$. Si, para todo esquema
$S'$ étale sobre~$S$ y todo par
\ifthenelse{\boolean{orig}}
{de objeto}
{de objetos}
$x_1,~y_1$ de~$(\Phi_1)_{S'}$, el haz $\SheafHom_{S'}(x_1,y_1)$ es
constructible, entonces, para todo par de objetos $x,~y$ de
$\Phi_{S'}$, lo mismo ocurre con~$\SheafHom_{S'}(x,y)$. Supongamos que
$\Phi_1$ es $1$-constructible (\Ref{XIII.0}) y que $\Phi_2$ es
constructible; entonces $\Phi$ es $1$-constructible.
\end{sublemme}

Para todo esquema $S'$ étale sobre~$S$ y todo objeto $x$ de
$\Phi_{S'}$, se tiene un monomorfismo
$$
\SheafAut_{S'}(x)\to\SheafAut_{S'}(p(x)).
$$
Resulta, pues, de SGA~4 IX~2.9 que, si $\Phi_1$ es constructible,
lo mismo ocurre con~$\Phi$, y la segunda afirmación del lema se
demuestra del mismo modo.

Supongamos ahora que $\Phi_1$ es $1$-constructible y que $\Phi_2$ es
constructible. El morfismo $p$ induce entre los haces de los subgerbes
maximales un morfismo
$$
\varphi\colon S\Phi\to S\Phi_1.
$$
Sea $G$ la imagen de~$S\Phi$ por~$\varphi$; según SGA~4 IX~2.9,
$G$ es un haz constructible. Puede encontrarse, por tanto, un haz
representable por un $S$-esquema étale de tipo finito~$T$ y un
epimorfismo
$$
a\colon T\to G
$$
(SGA~4 IX~2.7).
\marginpar{407}
Además, puede elegirse~$T$ de modo que la imagen~$y$ de la sección
identidad de~$T(T)$ por~$a$ esté definida por un objeto~$x_1$ de
$(\Phi_1)_T$ de la forma $x_1=p(x)$, donde $x\in\Phi_T$.

Basta demostrar que, para todo punto $s$ de~$S$, existe un abierto no
vacío~$U$ de~$\overline{\{ s\} }$ tal que $S\Phi|U$ sea localmente
constante y constructible. Sean $s\in S$, $\overline{s}$ un punto
geométrico sobre~$s$ y $\overline{y}_1,\dots,\overline{y}_n$ los
elementos de~$G_{\overline{s}}$. Por definición de~$T$, existen
morfismos $h_i\colon\overline{s}\to T$ tales que
$h_i^*(y)=\overline{y}_i$. Sea $S'$ el producto fibrado sobre~$S$ de
$n$ esquemas isomorfos a~$T$, sea $h\colon\overline{s}\to S'$ el
producto fibrado de los~$h_i$, y sean $y_i$ (\resp $x_i$) la imagen
inversa de~$y$ (\resp $x$) por la $i$-ésima proyección de~$S'$ sobre~$T$.
Sea $F_i$ el subhaz de~$S\Phi |S'$ que es la imagen inversa de~$y_i$, y demostremos
que los~$F_i$ son constructibles. El haz~$F_i$ es un cociente del haz
$F'_i$ tal que, para todo esquema~$S''$ étale sobre~$S'$, se tiene
\begin{multline*}
F'_i(S'') = \{\text{clases, salvo isomorfismo, de}
\ifthenelse{\boolean{orig}}
{\text{ objeto $z$ de $\Phi_{S''}$ provisto}}
{\text{ objetos $z$ de $\Phi_{S''}$ provistos}}\\
\ifthenelse{\boolean{orig}}
{\text{ de un isomorfismo $i\colon p(z)\simeq p(x_i|S'')$}.}
{\text{ de un isomorfismo $i\colon p(z)\simeq p(x_i|S'')$}\}.}
\end{multline*}
Basta demostrar que los~$F_i'$ son constructibles. Ahora bien, si se
pone $z_i=p_1p(x_i)$, $z_i'=p_2p(x_i)$, se tiene un monomorfismo
$$
\Psi_i\colon F'_i\to\SheafIsom_{S'}(z_i,z_i'),
$$
obtenido asociando, a todo esquema~$S''$ étale sobre~$S'$ y a todo
objeto~$z$ de~$\Phi_{S''}$ provisto de un isomorfismo
$i\colon p(z)\simeq p(x_i|S'')$, el isomorfismo de~$z_i$ en~$z_i'$
definido por la condición de que el diagrama
$$
\xymatrix@C=1.2cm{
p_1p(z) \ar[r]^-{p_1(i)}\ar[d]_-j& z_i\ar[d]\\ p_2p(z) \ar[r]^-{p_2(i)}& z'_i
}
$$
sea conmutativo ($j$ es el morfismo canónico asociado al diagrama
exacto).
\marginpar{408}

El morfismo~$\Psi_i$ es inyectivo, pues decir que dos objetos $z,~z'$,
provistos de isomorfismos $i\colon p(z)\isomto
p(x_i|S''),~i'\colon p(z')\isomto p(x_i|S'')$, definen el mismo
elemento de~$\SheafIsom_{S''}(z_i,z'_i)$ equivale a decir que
$p_1(i'^{-1}i)=p_2(i'^{-1}i)$, es decir, que $i'^{-1}i$ proviene de
un isomorfismo $z\to z'$. Como el haz~$F'_i$ es un subhaz de
$\SheafIsom_{S'}(z_i,z'_i)$, es constructible.

Como puede encontrarse un abierto no vacío~$U$ de~$\overline{\{s\}}$
tal que $y_1|U,\dots,y_n|U$ engendren~$G$, resulta del
lema~\Ref{XIII.6.1.2} que figura más abajo que $S\Phi |U$ es
constructible; por tanto, tras restringir~$U$, $S\Phi |U$ es localmente
constante y constructible.
